\chapter{Conclusions and Future Work}
\label{chap:conclusions}

\section{Conclusions}
\label{sec:conclusions}
This thesis asked whether a capable Vision--Language--Action model can be fine-tuned
and deployed on a single 12~GB consumer GPU to perform reliable, language-driven
grasping of everyday objects. The answer is a qualified yes. Working entirely in the
MuJoCo simulator, a compact SmolVLA policy was fine-tuned from 208 scripted-expert
demonstrations --- by freezing the pretrained vision--language backbone and training
only the action expert --- to pick a named grocery item from a shelf and place it in
an onboard basket. It does so reliably under the conditions it was trained on, and
markedly less reliably as soon as those conditions are perturbed; characterising that
boundary precisely, rather than reporting only the favourable figure, is among the
thesis's more useful outcomes.

The main outcomes are:
\begin{itemize}
    \item An 80\% task-completion rate (95\% Wilson CI 70--87, $n=80$) and 92.5\%
    grasp success for the selected checkpoint, drawn from 320 closed-loop rollouts
    across four checkpoints, with the deployed checkpoint chosen by task success
    rather than training loss (\refChap{chap:results}).
    \item A documented failure analysis that raised placement from 0\% ($n=12$) to 80\% by
    diagnosing near-miss behaviour and correcting observability and
    episode-boundary issues, rather than by scaling the model.
    \item Evidence that the policy is genuinely instruction-driven, together with the
    finding that this conditioning is implemented by \emph{recall rather than
    perception}: across 216 trials in which items were displaced from their trained
    positions, the policy reached the trained position every time and the actual
    position never. A light-weight view of the policy's own predicted action plan
    accompanies this.
    \item A measured bound on that behaviour: displacing objects by \SI{2.5}{cm} --- the
    same jitter present in the training data --- reduces task completion from 80\% to
    58.8\% (95\% CI 48--69).
    \item The finding --- corroborated by the literature and now by a same-task capacity
    control --- that evaluation quality and demonstration variety matter more than model
    scale in this data and hardware regime. A 52-million-parameter policy with no
    language input matches the 450-million-parameter VLA's manipulation quality on the
    item tested, locating the VLA's contribution in item selection.
\end{itemize}

Taken together, these show that a rigorously evaluated, compact VLA is a practical
choice under tight hardware constraints, that careful methodology is what turns an
initially failing policy into a working one, and --- equally --- that the same
methodology, applied without flinching to one's own headline claims, is what reveals
their limits. The evaluation designed in \refChap{chap:methodology} was ultimately
able to overturn a claim made earlier in this work; that is the strongest argument
available for building it that way.

More broadly, the thesis offers a template for working under constraint. When the
largest models are unavailable, progress comes not from waiting for more compute but
from three disciplines that the results here repeatedly reward: choosing the
smallest capable model and adapting only the part of it that must change (the action
mapping, with the perception prior kept frozen); investing in demonstration quality
and variety rather than volume; and measuring honestly, with task-grounded metrics,
adequate samples, confidence intervals and success-based model selection. None of
these is specific to SmolVLA or to a supermarket, and each transfers directly to the
navigation stage and to eventual hardware. The single most valuable outcome is
perhaps the failure analysis: it demonstrates that a well-instrumented diagnosis of
\emph{why} a policy fails is worth more than an untargeted increase in model size or
data, because it converts an opaque 0\% into a specific, fixable defect.

\section{Future work}
\label{sec:future}
The manipulation system studied here is the first half of a two-part vision of a
supermarket service robot. The clearest next steps are:
\begin{itemize}
    \item \textbf{A navigation VLA.} Extend the system with a second,
    language-guided policy that drives the mobile base between shelves in a store
    populated by moving, colliding shoppers. Given the safety and latency demands of
    a dynamic environment, a fast reactive policy (rather than a single large model)
    is appropriate, mirroring the dual-system argument of
    \refsec{sec:disc-ideal}.
    \item \textbf{An orchestrator for full delivery.} Combine navigation and
    manipulation so that the robot loops a complete shopping list --- navigate, pick,
    place --- and terminates on delivery, rather than performing one pick at a time.
    \item \textbf{Closing the multi-item gap.} Fine-tune on demonstrations with
    varied, non-empty basket contents to remove the distribution shift that reduces
    multi-item reliability (\refsec{sec:disc-limits}).
    \item \textbf{Toward hardware.} Domain randomisation and fine-tuning on real
    demonstrations (for example on a low-cost arm) to test sim-to-real transfer,
    which is out of scope for the present simulation study.
\end{itemize}

\subsection{The navigation stage in more detail}
\label{sec:future-nav}
The navigation stage is the most substantial extension and deserves elaboration,
because it raises problems the manipulation stage does not. Where manipulation is
quasi-static and the environment is fixed during a pick, navigation is dynamic:
shoppers move and may collide with the base, so the policy must be reactive and
operate under a real-time constraint. This changes the appropriate design. In
manipulation, executing a long 50-step action chunk before re-planning is beneficial,
because it yields smooth motion in a static scene; in navigation, committing to a long
horizon in a changing scene is dangerous, so a much shorter execution horizon (or
real-time re-planning) is required --- an instance of the non-monotonic
horizon effect established in the literature \parencite{wang2026vlalimits}. The methodology
would nonetheless mirror the manipulation pipeline: a scripted \emph{reactive
navigator} (for example a potential-field or sampling-based local planner that avoids
moving obstacles while heading to a goal shelf) would generate demonstrations of
language-goal-directed driving through a crowd; a navigation VLA would be fine-tuned on
these; and it would be evaluated closed-loop with success (reaching the goal without
collision) reported over many trials with confidence intervals, exactly as here.

Two further components complete the vision. First, an \emph{orchestrator} would loop
the full shopping list --- for each item, navigate to the relevant shelf, execute a
pick, and place --- terminating on delivery; this is a comparatively simple
sequencing layer over the two learned policies, and a first version already exists for
the manipulation-only case (\refsec{sec:results-orchestrator}). Second, the safety and
latency demands of driving near people argue for the dual-system architecture
discussed in \refsec{sec:disc-ideal}: a slower high-level module to interpret the list
and the scene, and a fast reactive policy for collision avoidance. The manipulation
results in this thesis provide the template --- controlled task, scripted expert,
frozen-backbone fine-tuning, and rigorous closed-loop evaluation --- that the
navigation stage would follow.

\subsection{Closing remark}
\label{sec:closing}
This thesis set out to test a single, practical proposition: that useful
language-driven manipulation does not require a frontier-scale model or a
frontier-scale budget, provided the method is disciplined. The evidence supports it.
A compact VLA, fine-tuned carefully and evaluated honestly on modest hardware, learned
to pick named everyday objects and place them reliably, and its one categorical
failure was resolved not by scaling but by understanding it. That is an encouraging
result for anyone seeking to build capable robots under real-world constraints.
